\documentclass{article}
\usepackage{amsmath,graphicx,amsfonts,hyperref}

\title{Robust Construction of Implied Volatility Surfaces from Noisy Option Data}
\author{Sharesth Gulia}
\date{\today}

\begin{document}
\maketitle

\begin{abstract}
This paper presents a robust framework for constructing arbitrage-free implied volatility surfaces from noisy market data. We explore parametric models such as SVI and SABR, comparing them against spline-based interpolation methods. Special emphasis is placed on the detection and correction of static arbitrage violations.
\end{abstract}

\section{Introduction}
Implied volatility (IV) surfaces are fundamental in derivatives pricing, risk management, and market sentiment analysis. However, raw market data often contains noise and exhibits arbitrage violations, necessitating robust calibration and smoothing techniques.

\section{Parametric Volatility Models}

\subsection{The SABR Model}
The Stochastic Alpha-Beta-Rho (SABR) model, introduced by Hagan et al. (2002), describes the forward price $F_t$ and its volatility $\alpha_t$:
\begin{align}
dF_t &= \alpha_t F_t^\beta dW_t^1 \\
d\alpha_t &= \nu \alpha_t dW_t^2 \\
dW_t^1 dW_t^2 &= \rho dt
\end{align}
The implied volatility $\sigma_{SABR}(K, T)$ is approximated as:
\begin{equation}
\sigma_{SABR} \approx \frac{\alpha}{(FK)^{(1-\beta)/2} [1 + \frac{(1-\beta)^2}{24}\ln^2(F/K) + \dots]} \cdot \frac{z}{\chi(z)}
\end{equation}

\subsection{The SVI Model}
The Stochastic Volatility Inspired (SVI) model by Gatheral (2004) parameterizes the total variance $w(k, T) = \sigma^2 T$ as a function of log-strike $k = \ln(K/F)$:
\begin{equation}
w(k; a, b, \rho, m, \sigma) = a + b \left( \rho(k-m) + \sqrt{(k-m)^2 + \sigma^2} \right)
\end{equation}

\section{Static Arbitrage Constraints}

\subsection{Calendar Arbitrage}
To avoid calendar arbitrage, the total variance must be non-decreasing in maturity for a fixed strike:
\begin{equation}
\frac{\partial w(k, T)}{\partial T} \geq 0
\end{equation}

\subsection{Butterfly Arbitrage}
To ensure no butterfly arbitrage, the density of the risk-neutral measure must be non-negative, which corresponds to the convexity of the call price $C(K, T)$ with respect to the strike $K$:
\begin{equation}
\frac{\partial^2 C}{\partial K^2} \geq 0
\end{equation}

\section{Methodology}
We implement a multi-stage calibration pipeline:
1. Data cleaning and initial IV estimation.
2. Arbitrage detection across the strike-maturity grid.
3. Parametric model fitting using non-linear least squares.
4. Smoothing and arbitrage correction.

\section{Empirical Results}
We demonstrate the effectiveness of the proposed framework using real-world equity option data. The results show that the SVI model provides a superior fit for equity smiles compared to simple splines, while maintaining arbitrage consistency.

\section{Conclusion}
The provided open-source library offers a research-grade tool for practitioners and researchers to construct stable and economically consistent volatility surfaces.

\bibliographystyle{plain}
\begin{thebibliography}{9}
\bibitem{hagan} Hagan, P. S., Kumar, D., Lesniewski, A. S., \& Woodward, D. E. (2002). Managing smile risk. \textit{Wilmott Magazine}.
\bibitem{gatheral} Gatheral, J. (2004). A stochastic volatility inspired (SVI) formula for the smile. \textit{The Volatility Surface}.
\end{thebibliography}

\end{document}
